%----------------------------------------------------------------------------------------
%	PACKAGES AND THEMES
%----------------------------------------------------------------------------------------
\documentclass{beamer}
% Load a bunch of useful packages:
\usepackage{amssymb,amsmath,amsfonts,mathtools} % useful math fonts and symbols
\usepackage{geometry} % allows changing margins and sizes of stuff
\usepackage{hyperref} % allows referencing of lines of text, urls, figures etc
\usepackage{natbib} % allows citation referencing from a .bib file
\bibliographystyle{apalike} % American Psychological Association style guide
\usepackage{graphicx} % to include graphics from the figures folder with helpful formating options
%---------------------------------------------%
\usepackage{color} % custom color definitions:%
    %UOregon colors:                          %
    \definecolor{UOGreen}{RGB}{18, 71, 52}    %
    \definecolor{UOYellow}{RGB}{254, 225, 35} %
    %Secondary official colors:               %
    \definecolor{LegacyGreen}{HTML}{104735}   %
    \definecolor{GrassGreen}{HTML}{489D46}    %
    \definecolor{LimeGreen}{HTML}{8ABB40}     %
    \definecolor{Chartreuse}{HTML}{E2E11B}    %
    \definecolor{Berry}{HTML}{8D1D58}         %
    \definecolor{DarkBlue}{HTML}{004F6E}      %
    \definecolor{LightBlue}{HTML}{00A5B5}     %
%----------------------------------------------
\usepackage{setspace} % to specify single, double or one-half spacing of text
\usepackage{indentfirst} % indent first line of a text paragraph
\usepackage{multicol} %multipage ability
\usepackage{ulem} % \ul (underline) command which will break over line ends
\usepackage{amsthm} % allows standarized theorem commands
\setbeamertemplate{theorems}[numbered]
\theoremstyle{plain}
\newtheorem{assume}{Assumption}
\usepackage{breqn} %automatic line breaking
\normalem

\usetheme{Frankfurt}
%UOregon colors:
\definecolor{UOGreen}{RGB}{18, 71, 52}
\definecolor{UOYellow}{RGB}{254, 225, 35}
%Secondary official colors:
\definecolor{LegacyGreen}{HTML}{104735}
\definecolor{GrassGreen}{HTML}{489D46}
\definecolor{LimeGreen}{HTML}{8ABB40}
\definecolor{Chartreuse}{HTML}{E2E11B}
\definecolor{Berry}{HTML}{8D1D58}
\definecolor{DarkBlue}{HTML}{004F6E}
\definecolor{LightBlue}{HTML}{00A5B5}
\usefonttheme{professionalfonts}
\fontfamily{ppl}\selectfont
\usecolortheme[named=LegacyGreen]{structure}
\setbeamertemplate{footline}[frame number]
\setbeamertemplate{headline}{} %Removes Section Index on top of slide
\hypersetup{
    colorlinks=true,
    linkcolor=GrassGreen,
    filecolor=magenta,      
    urlcolor={LightBlue},
    pdftitle={Overleaf Example},
    pdfpagemode=FullScreen
}

%----------------------------------------------------------------------------------------
%	TITLE PAGE
%----------------------------------------------------------------------------------------

% The title
\title{Forced Removal or Economic Mobility?}
\subtitle{Wage Effects of Japanese Internment}
\author{Dante Yasui}
\institute{Presented in Dept. of Economics Grad Group}
\date{\today}
\titlegraphic{
\includegraphics[scale=.4]{figures/UOSignature-356.png}}

\begin{document}

\maketitle

%---------------------------------------------------------------------------
\section{Introduction}
%---------------------------------------------------------------------------

\begin{frame}{Japanese Internment as Natural Experiment}
    \begin{figure}
        \centering
        \includegraphics[width=0.7\linewidth]{figures/PioCrtInternmentMap.jpg}
    \end{figure}
    {\footnotesize \href{https://pioneercourthouse.org/gallery-Japanese-American-Internment.html}{Source}}
\end{frame}

\begin{frame}{Pre-Internment Japanese Distribution}\label{ctmap}
    \begin{figure}
        \centering
        \includegraphics[width=.85\linewidth]{figures/county_JAmap.png}
        \label{fig:county40map}
    \end{figure}
\end{frame}

%---------------------------------------------------------------------------
\section{Research Question}
%---------------------------------------------------------------------------

\begin{frame}{Research Question}\label{rq_slide}
    \begin{block}{Research Question:}
        \textbf{How were internee labor market outcomes affected after their release?}
    \end{block}
    \begin{itemize}
        \item WRA camps and resettlement lead to shift in Japanese Population \hyperlink{areaplot}{\beamergotobutton{Plot}}
        \item Variation in camp areas by local median incomes \hyperlink{stincomemap}{\beamergotobutton{Map}}
        \item Camps assigned by geography and capacity; not economic characteristics \hyperlink{campplcmap}{\beamergotobutton{Map}}
    \end{itemize}
\end{frame}

%---------------------------------------------------------------------------
\section{Significance and Contribution}
%---------------------------------------------------------------------------

\begin{frame}{Related Literature}
    \begin{block}{Effects on Japanese American Internees:}
        \begin{itemize}
            \item Internment caused labor market withdrawl leading to reduction of working-age male earnings by 9-13\% in 1970 \citep{chin_longrun_2005}
            \item Internment may have lead to higher long-term wages due to re-optimization and camp diversity (Arellano-Bover 2022).
            \item Camp assignment induced resettlement and those moved to richer areas did better then those moved to poorer areas in 1980s \citep{shoag_causal_2016}
        \end{itemize}
    \end{block}
\end{frame}

%---------------------------------------------------------------------------
\section{Methods}
%---------------------------------------------------------------------------

\begin{frame}{Data}
    \textbf{Sources:} 
    \begin{itemize}
        \item 1940 and 1950 full count census via IPUMS \citep{ruggles_steven_ipums_2020}
          \begin{itemize}
            \item I use IPUMS' linked data: the Multigenerational Longitudinal Panel (MLP)
           \end{itemize}
        \item Wartime Relocation Authority (WRA) camp records via National Archives
    \end{itemize}
    \textbf{Variables:}
    \begin{itemize}
    \item Outcomes: wages, employment, occupation
    \item Time-Invariant: birth year, birthplace, sex, etc
    \item Location: either census address or 1942 address prior to EO9066 evacuations
    \item Treatment: predicted internment status
    \end{itemize}
\end{frame}

    \begin{frame}{Methods}
      \begin{block}{Predicted Internment Status}
        Using known proportions of demographic categories within WRA and census,
        we can estimate the probability of internment using Bayes Rule:
        $$
        Pr(I_i=1|Z_i,c_i^{40}) = \frac{Pr(z_i,c_i^{40}|I_i=1) \cdot Pr(I_i=1)}{Pr(z_i,c_i^{40})}
        $$
        \begin{itemize}
        \item $Pr(z_i,c_i^{40}|I_i=1)$: proportion of internees in each county-demographic
        \item $Pr(I_i=1)$: fraction of internees over 1940 Japanese American population
        \item $Pr(z_i,s_i^{40})$: proportion of Japanese Americans in each county-demographic in 1940 census
        \end{itemize}
      \end{block}
    \end{frame}

    \begin{frame}
      \begin{table}
        \begin{tabular}{lrrrr}
          \textbf{YEAR:} & \multicolumn{2}{c}{1940} & \multicolumn{2}{c}{1950} \\
          \underline{variable} & Chinese & Japanese & Chinese & Japanese \\
          $\hat{E}[\text{internment}]$: & 0 &  0.66 & 0 &  0.64 \\
          Percent Female: &  0.33 &  0.45 &  0.34 &  0.45 \\
          Average BirthYear: & 1912 & 1908 & 1913 & 1908 \\
          Employed: & 0.37 & 0.42 &  0.50 & 0.52 \\
          Some College: & 0.05 & 0.05 & 0.036 & 0.03 \\
        \end{tabular}
      \end{table}
    \end{frame}

    \begin{frame}{Summary Statistics}
      \begin{table}[!t]                                        
        \fontsize{12.0pt}{14.4pt}\selectfont                     
        \begin{tabular*}{\linewidth}{@{\extracolsep{\fill}}lcc}  
          \textbf{Characteristic} & \textbf{4}\\                   
          N = 6,448\textsuperscript{\textit{1}} & \textbf{5}\\     
          N = 4,631\textsuperscript{\textit{1}} \\                 
          \midrule\addlinespace[2.5pt]                             
          pr\_intern & 0.00 (0.00, 0.00) & 0.84 (0.40, 1.00) \\    
          female & 2,157 (33\%) & 2,101 (45\%) \\                  
          BIRTHYR & 1,912 (1,899, 1,930) & 1,904 (1,892, 1,927) \\ 
          employed & 2,388 (37\%) & 1,947 (42\%) \\                
          college & 340 (5.3\%) & 244 (5.3\%) \\                   
          foreign & 181 (2.8\%) & 149 (3.2\%) \\                   
        \end{tabular*}                                           
        \begin{minipage}{\linewidth}                             
          \textsuperscript{\textit{1}}Median (Q1, Q3); n (\%)\\    
        \end{minipage}                                           
      \end{table}                                              
    \end{frame}

    \begin{frame}
      \begin{figure}
        \centering
        \includegraphics[width=0.7\linewidth]{figures/did_plot.png}
      \end{figure}
    \end{frame}
    
    \begin{frame}{Methods}\label{methods}
      \begin{block}{Difference-in-Differences \hyperlink{DiDmethod}{\beamergotobutton{Diff-in-Diff}}}
        Assuming that Chinese and Japanese Americans would have experienced similar wage trends without internment:
        $$y_{it} = \alpha_t + X_{it}' \gamma + \delta \hat{E}[I_i|Z_i,c_i^t] + \beta (\hat{E}[I_i|Z_i,c_i^t] \times Post_t) + \epsilon_{it}$$
      \end{block}
        \begin{itemize}
        \item $y_{it}$: real wages earned by individual $i$ in year $t$
        \item $\hat{E}[I_i|Z_i,c_i^t]$: Predicted internhment status
        \item $Post_t$: treatment has an effect in 1950, but not 1940
        \item $X_{it}$: Time-variant controls (educ, marst, etc)
        \end{itemize}
        $\beta$ : DiD estimate between internees and non-internees, pre vs post internment
    \end{frame}

    \begin{frame}{Diff-in-Diff}
      \begin{block}{Identifying assumption:}
        The only reason an internee would resettle and stay in destination $k$ is through the internment and resettlement mechanism
      \end{block} 
      \begin{itemize}
      \item Possibly violated if internees did have some degree of choice and ability to move between release and 1950 census
      \end{itemize}
    \end{frame}

    \begin{frame}
      \begingroup
      \centering
      \begin{tabular}{lcccc}
        \tabularnewline \midrule \midrule
        Dependent Variable: & \multicolumn{4}{c}{INCWAGE}\\
        Model:                & (1)            & (2)            & (3)            & (4)\\  
        \midrule
        \emph{Variables}\\
        Constant              & 334.0$^{***}$  & -1,591.9       &                &   \\   
        & (9.644)        & (1,030.1)      &                &   \\   
        did                   & 170.9$^{***}$  & 189.8$^{***}$  & 36.45          & 62.72\\   
        & (33.62)        & (32.36)        & (40.87)        & (33.43)\\   
        pr\_intern            & -121.7$^{***}$ & -62.10$^{***}$ & -39.38$^{***}$ & -27.95$^{***}$\\   
        & (18.54)        & (18.17)        & (7.830)        & (6.908)\\   
        as\_factor(YEAR)1950  & 174.9$^{***}$  & 203.0$^{***}$  & 295.3$^{***}$  & 321.8$^{***}$\\   
        & (17.99)        & (17.99)        & (40.59)        & (35.45)\\   
        female                &                & -347.9$^{***}$ &                & -287.3$^{***}$\\   
        &                & (13.65)        &                & (7.192)\\   
        BIRTHYR               &                & 0.9884$^{*}$   &                & 0.8518$^{***}$\\   
        &                & (0.5377)       &                & (0.2067)\\   
        college               &                & 233.5$^{***}$  &                & 129.3$^{***}$\\   
        &                & (22.34)        &                & (11.67)\\   
        married               &                & 174.6$^{***}$  &                & 173.2$^{***}$\\   
        &                & (17.64)        &                & (7.610)\\   
        \midrule
        \emph{Fixed-effects}\\
        as\_factor(state)     &                &                & Yes            & Yes\\  
        \midrule
        \emph{Fit statistics}\\
        Observations          & 10,309         & 10,309         & 4,720          & 4,720\\  
        R$^2$                 & 0.02777        & 0.10337        & 0.05570        & 0.10621\\  
        Within R$^2$          &                &                & 0.05040        & 0.10120\\  
        \midrule \midrule
        \multicolumn{5}{l}{\emph{Signif. Codes: ***: 0.01, **: 0.05, *: 0.1}}\\
      \end{tabular}
      \par\endgroup
    \end{frame}

    \begin{frame}
      \begingroup
      \centering
      \begin{tabular}{lccc}
        \tabularnewline \midrule \midrule
        Dependent Variables:  & employed        & OCCSCORE        & college\\  
        Model:                & (1)             & (2)             & (3)\\  
        \midrule
        \emph{Variables}\\
        Constant              & 8.696$^{***}$   & 162.2$^{***}$   & $-3.96\times 10^{-14}$\\    
        & (0.4347)        & (12.58)         & ($1.19\times 10^{-13}$)\\    
        did                   & -0.0739$^{***}$ & -1.977$^{***}$  & $-3.38\times 10^{-17}$\\    
        & (0.0138)        & (0.3988)        & ($3.78\times 10^{-15}$)\\    
        pr\_intern            & 0.0779$^{***}$  & -1.797$^{***}$  & $-1.48\times 10^{-16}$\\    
        & (0.0098)        & (0.2840)        & ($2.69\times 10^{-15}$)\\    
        as\_factor(YEAR)1950  & 0.1396$^{***}$  & 2.335$^{***}$   & $-2.11\times 10^{-17}$\\    
        & (0.0068)        & (0.1968)        & ($1.87\times 10^{-15}$)\\    
        female                & -0.3449$^{***}$ & -9.011$^{***}$  & $2.84\times 10^{-16}$\\    
        & (0.0059)        & (0.1720)        & ($1.63\times 10^{-15}$)\\    
        BIRTHYR               & -0.0044$^{***}$ & -0.0802$^{***}$ & $8.12\times 10^{-17}$\\    
        & (0.0002)        & (0.0065)        & ($6.2\times 10^{-17}$)\\    
        college               & 0.1216$^{***}$  & 8.288$^{***}$   & 1$^{***}$\\   
        & (0.0139)        & (0.4033)        & ($3.83\times 10^{-15}$)\\    
        married               & 0.2354$^{***}$  & 7.285$^{***}$   & $2.87\times 10^{-16}$\\    
        & (0.0083)        & (0.2396)        & ($2.27\times 10^{-15}$)\\    
        \midrule
        \emph{Fit statistics}\\
        Observations          & 22,158          & 22,158          & 22,158\\  
        R$^2$                 & 0.28101         & 0.24857         & 1\\  
        Adjusted R$^2$        & 0.28078         & 0.24833         & 1\\  
        \midrule \midrule
        \multicolumn{4}{l}{\emph{IID standard-errors in parentheses}}\\
        \multicolumn{4}{l}{\emph{Signif. Codes: ***: 0.01, **: 0.05, *: 0.1}}\\
      \end{tabular}
      \par\endgroup
    \end{frame}
    
    \section{Conclusion}

    \begin{frame}[plain]{}
      \tiny
      \bibliography{InternRefs}
      \addtocounter{framenumber}{-1}
    \end{frame}

    \section{maps}

    \begin{frame}[plain]\label{wrapmap}
      \begin{figure}
        \centering
        \includegraphics[width=.95\linewidth]{figures/wrap_maps.png}
      \end{figure}
      \addtocounter{framenumber}{-1}
    \end{frame}

    \begin{frame}[plain]\label{areaplot}
      \begin{figure}
        \caption{Regional Trends in Japanese Population Concentration \hyperlink{rq_slide}{\beamerreturnbutton{Back}}}
        \centering
        \includegraphics[width=.65\linewidth]{figures/shareareaplot.png}
        \label{fig:areaplot}
      \end{figure}
      \addtocounter{framenumber}{-1}
    \end{frame}

    \begin{frame}[plain]\label{stincomemap}
      \begin{figure}
        \centering
        \includegraphics[height=.6\linewidth]{figures/st_income_map.png}
      \end{figure}
      \hyperlink{rq_slide}{\beamerreturnbutton{Back}}
      \addtocounter{framenumber}{-1}
    \end{frame}
    
    \begin{frame}[plain]\label{campplcmap}
      \begin{figure}
        \centering
        \includegraphics[height=.6\linewidth]{figures/RDmap.png}
      \end{figure}
      \hyperlink{rq_slide}{\beamerreturnbutton{Back}}
                {\footnotesize \href{http://resource.nlm.nih.gov/01130040R}{Source}}
                \addtocounter{framenumber}{-1}
    \end{frame}

\end{document}

